\section{The classical distribution relations and their arithmetic consequences}
\label{sec:kummer-distribution}

The multiplicative relations among Ihara's Kummer cocycles already hold
integrally on the full profinite Grothendieck--Teichm\"uller group.
They follow from Enriquez's distribution theorem, which is used here as
published background.

For the cyclic-cover subgroup $H_n\subset\widehat F_2$, define
\[
 d_n:H_n\longrightarrow\widehat F_2,\qquad
 x^n\longmapsto x,\quad y_0\longmapsto y,\quad
 y_j\longmapsto1\quad(0<j<n).
\]
Let $c_n=\operatorname{exp}_y\circ d_n$; thus $c_n=\ell_n$ in
our Schreier coordinates. Enriquez proves that, for every
$g=(\lambda,f)\in\widehat{\mathrm{GT}}$ and every $n\geq1$,
\begin{equation}
 d_n(f)=y^{c_n(f)}f.
 \label{eq:full-kummer-distribution}
\end{equation}
This is equation (50) and Theorem 8.1 of \cite{Enriquez2007}; in
particular his subgroup $\mathrm{GTK}$ is the full profinite group.
The map is the filling map on the cyclic cover, so this is an identity
in the full free profinite group, beyond its abelian or congruence images.

\begin{proposition}\label{prop:kummer-multiplicativity}
For positive integers $a,b$ and every $g\in\widehat{\mathrm{GT}}$,
\[
 \rho_{ab}^I(g)=\rho_a^I(g)+\rho_b^I(g).
\]
Consequently
\[
 \rho_n^I=\sum_{q\mid n}v_q(n)\rho_q^I,
 \qquad
 \frac{\rho_{ab}^I-\rho_a^I-\rho_b^I}{2}\bmod2=0.
\]
The last expression therefore supplies no nonzero character in the
restriction kernel.
\end{proposition}
\begin{proof}
On $H_{ab}$ one has $d_b(H_{ab})\subset H_a$ and
$d_a d_b=d_{ab}$. Indeed $x^{ab}$ maps successively to $x^a,x$;
for $y_j=x^jyx^{-j}$, the first map kills it unless $b\mid j$,
and otherwise sends it to $y_{j/b}$. The second map retains only
$j=0$, for $0\leq j<ab$.
Applying \eqref{eq:full-kummer-distribution} twice now gives
\[
 y^{c_{ab}(f)}f=d_{ab}(f)
   =d_a\bigl(y^{c_b(f)}f\bigr)
   =y^{c_b(f)+c_a(f)}f.
\]
The procyclic subgroup generated by $y$ is isomorphic to
$\widehat{\mathbb Z}$, so $c_{ab}(f)=c_a(f)+c_b(f)$.
This is also the observation in \cite[Remark 8.6]{Enriquez2007}.
Since $\rho_n^I=-c_n(f)$, the asserted formulas follow, with the
prime factorization formula obtained by iteration.
\end{proof}

The cocycle normalization agrees with the one used for the modular
frame: \cite[\S2.2, footnote on p.~100]{Furusho2003} identifies the
Nakamura--Schneps cocycle $\rho_2$ with $-\Psi_2^{(0)}$, and
\cite[Lemma 7(a) and Proposition 11]{H17} identifies that cocycle
with the pure-frame definition. Hence $\rho_2^I=\rho_2$ on
$\widehat{\mathrm{GT}}$. The negative-index Schreier generators in
Furusho's convention give the same coefficient $c_n$: for $0<j<n$,
$x^{-j}yx^j$ is conjugate within $H_n$ to $y_{n-j}$, so both
coefficient maps vanish on it.

\begin{corollary}\label{cor:dyadic-anderson}
Every element of $\widehat{\mathrm{GT}}$ satisfies all dyadic
Anderson multiplication relations $(A_{2^j})$, $j\geq1$, in the
notation of \cite[Definition 2.3]{Furusho2003}.
\end{corollary}
\begin{proof}
Relation (IV) holds on the group by \cite[Theorem 1]{H17}, and
$(K_2)$ holds by \eqref{eq:full-kummer-distribution}.
\cite[Theorem 3.1]{Furusho2003} therefore gives the duplication
relation; its dyadic iterates follow from that paper's Lemma 3.2.
Equivalently, combine
$\mathrm{GTK}=\widehat{\mathrm{GT}}$,
$\mathrm{I}\Gamma=\widehat{\mathrm{GT}}$
\cite[Corollary 21]{H17}, and
\cite[Theorem 3.4]{Furusho2003}.
\end{proof}

These are consequences of the cited distribution and multiplication
theorems. They do not determine the full restriction kernel: they show
that the displayed multiplicative defects vanish and impose the
specified dyadic identities.
